% Manuscript Figure 5, on its own, with its legend.
%
% The figure PDF is a validation output and lives outside this repository, so this document
% refers to it by bare filename and must be compiled from the directory that holds it:
%
%   cp manuscript/figure5.tex D:/research/nasage_validation/results/mvri_audit_trim/figures/
%   cd D:/research/nasage_validation/results/mvri_audit_trim/figures
%   tectonic -X compile figure5.tex
%
\documentclass[11pt]{article}
\usepackage[a4paper,margin=20mm]{geometry}
\usepackage{graphicx}
\usepackage[T1]{fontenc}
\usepackage{lmodern}
\usepackage{microtype}
\usepackage{parskip}
\setlength{\parindent}{0pt}
\pagestyle{empty}

\newcommand{\figpage}[1]{%
  \thispagestyle{empty}%
  \vspace*{\fill}%
  \begin{center}\includegraphics[width=\textwidth,height=\textheight,keepaspectratio]{#1}\end{center}%
  \vspace*{\fill}%
  \clearpage}

\begin{document}

\figpage{figure5_maturation.pdf}

{\large\textbf{Figure 5}\par}\vspace{2pt}
{\textbf{Maturation of the mucosal interface across the lifespan.}\par}\vspace{8pt}

\small

Five candidate axes of the mucosal interface --- mucus permeability, antigen presentation,
innate sensing, IgA support and cytokine tone --- were scored in 148,978 nasal cells from
52 donors spanning seven age bins, using the gene lists and cell subsets defined for each
axis. Cells were labelled with \texttt{Cell\_type\_prediction\_level1/2}, which covers every
cell in the object. Two axes are absent from this figure: innate sensing and cytokine tone
fail a measurability gate, with only 45\% and 54\% of their genes detected in at least 5\% of
the cells they would be scored in, so a flat trajectory for either would mean \emph{not
measured} rather than \emph{not different}. Within each remaining axis, genes below that same
5\% detection floor are excluded from the score, so the axes shown are built from 13 of 15,
11 of 17 and 3 of 4 declared genes respectively.

\textbf{a}, What each measurable axis is built from and where it is read out. Node and ribbon
heights are flows: an axis is as tall as the number of genes that survive the detection
floor, and that height is divided among cell types in proportion to the cells the axis is
scored in. The composition is itself informative --- all three axes are read out
predominantly in Goblet and Secretory cells, and the antigen-presentation axis draws under
2\% of its cells from the myeloid compartment. Cell nodes are drawn with a minimum height so
that small compartments remain visible; the counts beside each node are exact.

\textbf{b}, Donor-level trajectories. Each point is one donor, coloured by age bin, and the
bar is the bin median; $n$ per bin is printed on the axis. Scores are donor means of the
per-gene $z$ within each axis's own cell subset. The headline statistic is the empirical
$P$ against 500 expression- and size-matched random gene sets, scored and tested identically
at donor level: antigen presentation $P < 0.002$ (0 of 500 null sets reached the observed
statistic) and IgA support $P = 0.002$ both clear it, while mucus permeability does not
($P = 0.066$). Benjamini--Hochberg values across the five axes tested are given in
parentheses but should not be read as the result: because the three source datasets in this
cohort are confounded with age bin, 52--75\% of \emph{random} matched gene sets reach a
nominal donor-level $P < 0.05$, so the nominal and BH values are not calibrated here and the
matched-null comparison is the only valid test.

\textbf{c}, Spearman correlation of donor score with age-bin rank within two age windows,
with 95\% bootstrap confidence intervals. The maturation is concentrated between birth and
young childhood, where all three axes rise steeply (mucus permeability $\rho = +0.56$,
$P = 0.003$; antigen presentation $\rho = +0.81$, $P < 0.001$; IgA support $\rho = +0.80$,
$P < 0.001$; $n = 26$ donors). After childhood, antigen presentation and IgA support decline
($\rho = -0.55$, $P = 0.003$ and $\rho = -0.42$, $P = 0.035$) while mucus permeability does
not change ($\rho = +0.07$, $P = 0.75$).

\textbf{d}, Mean score by cell type within each axis's own subset. Dot colour is the mean
per-gene $z$ in that group and dot size the number of cells.

\textbf{e}, The three measurable axes combined into a single score, using the declared axis
signs renormalised over the surviving axes. The combined score rises from the neonatal period
to a childhood maximum and falls thereafter. The three axes are not three independent
measurements: they give 2.0 effective tests by the Li and Ji criterion, with the first
principal component carrying 73\% of the variance and pairwise Spearman correlations of
$+0.53$ to $+0.70$.

\vspace{6pt}
\textbf{Interpretation.} The mucosal interface is not static across the lifespan, and the
part of its trajectory that the data support most strongly is its construction rather than
its decline: between birth and young childhood, antigen presentation and IgA support rise
steeply and both beat expression-matched random gene sets. This is consistent with Figure 4,
where antigen-presentation and interferon programmes peak in childhood, and it places that
peak in a broader interface context by adding the secretory and IgA arms.

\textbf{Limitations.} Three bear directly on how far this figure can be read. First, the
decline after childhood cannot be separated from batch: the Elderly bin is four-fifths drawn
from one source dataset and the Adult bin entirely from another, and 10x chemistry tracks
that split, so the downward arm of panels c and e is confounded by design. Second, sequencing
depth correlates with every axis score ($\rho = +0.52$ to $+0.55$) and accounts for part of
the coherence between axes (mean pairwise $\rho$ falls from $+0.59$ to $+0.40$ once depth is
removed); it is not, however, confounded with age ($\rho = +0.10$, $P = 0.49$), and the age
$P$-values are unchanged by adjusting for it. Third, no vaccine or infection outcome exists in
this cohort, so these axes describe the maturation of the interface and cannot speak to
mucosal vaccine responsiveness.

\end{document}
